% p360_1; source PDF page 360; manually verified editable table.
\begin{tikzpicture}[x=1cm,y=-1cm,inner sep=0pt,outer sep=0pt]
\path[use as bounding box] (0,0) rectangle (16.000,10.930);
\draw[line width=.4pt] (0.000,0.000) rectangle (2.600,0.805);
\node[align=center,text width=2.340cm,font={\linespread{1}\fontsize{10}{13.50}\selectfont\bfseries}] at (1.300,0.403) {妨害公务罪};
\draw[line width=.4pt] (2.600,0.000) rectangle (5.700,0.805);
\node[align=center,text width=2.840cm,font={\linespread{1}\fontsize{10}{13.50}\selectfont\bfseries}] at (4.150,0.403) {内容};
\draw[line width=.4pt] (5.700,0.000) rectangle (16.000,0.805);
\node[align=center,text width=10.040cm,font={\linespread{1}\fontsize{10}{13.50}\selectfont\bfseries}] at (10.850,0.403) {范例};
\draw[line width=.4pt] (0.000,0.805) rectangle (2.600,2.610);
\node[align=center,text width=2.340cm,font={\linespread{1}\fontsize{10}{13.50}\selectfont}] at (1.300,1.708) {妨害的对象};
\draw[line width=.4pt] (2.600,0.805) rectangle (5.700,2.610);
\node[align=center,text width=2.840cm,font={\linespread{1}\fontsize{10}{13.50}\selectfont}] at (4.150,1.708) {正在执行的公务};
\draw[line width=.4pt] (5.700,0.805) rectangle (16.000,2.610);
\node[align=left,text width=10.040cm,font={\linespread{1}\fontsize{10}{13.50}\selectfont}] at (10.850,1.708) {警察调解完甲乙的纠纷，甲乙签字后，甲心里又不平，用脚踹警察，致其轻伤。甲不构成妨害公务罪，而构成故意伤害罪};
\draw[line width=.4pt] (0.000,2.610) rectangle (2.600,8.320);
\node[align=center,text width=2.340cm,font={\linespread{1}\fontsize{10}{13.50}\selectfont}] at (1.300,5.465) {认识错误};
\draw[line width=.4pt] (2.600,2.610) rectangle (5.700,5.415);
\node[align=center,text width=2.840cm,font={\linespread{1}\fontsize{10}{13.50}\selectfont}] at (4.150,4.013) {故意伤害罪};
\draw[line width=.4pt] (5.700,2.610) rectangle (16.000,5.415);
\node[align=left,text width=10.040cm,font={\linespread{1}\fontsize{10}{13.50}\selectfont}] at (10.850,4.013) {甲看到朋友乙与丙打架，便上前帮忙，将丙打成轻伤，实际上丙是便衣警察，正在抓捕乙。（1）妨害公务罪：客观要件（$\checkmark$），主观要件（$\times$）。（2）故意伤害罪：客观要件（$\checkmark$），主观要件（$\checkmark$），定故意伤害罪};
\draw[line width=.4pt] (2.600,5.415) rectangle (5.700,8.320);
\node[align=center,text width=2.840cm,font={\linespread{1}\fontsize{10}{13.50}\selectfont}] at (4.150,6.867) {假想防卫};
\draw[line width=.4pt] (5.700,5.415) rectangle (16.000,8.320);
\node[align=left,text width=10.040cm,font={\linespread{1}\fontsize{10}{13.50}\selectfont}] at (10.850,6.867) {甲误以为丙在追杀乙，便将丙打成重伤，实际上丙是便衣警察，正在追捕乙。（1）妨害公务罪：客观要件（$\checkmark$），主观要件（$\times$）。（2）故意伤害罪：客观要件（$\checkmark$），主观要件（$\times$）。如果有过失，则是过失犯罪；如果没有过失，则无罪（2013 年主观题）};
\draw[line width=.4pt] (0.000,8.320) rectangle (2.600,10.930);
\node[align=center,text width=2.340cm,font={\linespread{1}\fontsize{10}{13.50}\selectfont}] at (1.300,9.625) {罪数总结};
\draw[line width=.4pt] (2.600,8.320) rectangle (16.000,9.225);
\node[align=left,text width=13.140cm,font={\linespread{1}\fontsize{10}{13.50}\selectfont}] at (9.300,8.773) {原则上，实施某种犯罪，又实施妨害公务罪，应数罪并罚};
\draw[line width=.4pt] (2.600,9.225) rectangle (16.000,10.930);
\node[align=left,text width=13.140cm,font={\linespread{1}\fontsize{10}{13.50}\selectfont}] at (9.300,10.078) {例外的，结合犯：前罪+妨害公务罪=前罪，加重处罚。这种前罪只有：（1）走私、贩卖、运输、制造毒品犯罪。（2）组织、运送他人偷越国（边）境罪};
\end{tikzpicture}
